\section{Coherent formal modules}
\label{section-coherent-formal}

\noindent
As we do not yet have the theory of formal schemes to our disposal, we
develop a bit of language that replaces the notion of a ``coherent module on a
Noetherian adic formal scheme''.

\medskip\noindent
Let $X$ be a Noetherian scheme and let $\mathcal{I} \subset \mathcal{O}_X$
be a quasi-coherent sheaf of ideals. We will consider inverse
systems $(\mathcal{F}_n)$ of coherent $\mathcal{O}_X$-modules such that
\begin{enumerate}
\item $\mathcal{F}_n$ is annihilated by $\mathcal{I}^n$, and
\item the transition maps induce isomorphisms
$\mathcal{F}_{n + 1}/\mathcal{I}^n\mathcal{F}_{n + 1} \to \mathcal{F}_n$.
\end{enumerate}
A morphism of such inverse systems is defined as usual.
Let us denote the category of these inverse systems with
$\textit{Coh}(X, \mathcal{I})$. We are going to proceed by proving
a bunch of lemmas about objects in this category. In fact, most
of the lemmas that follow are straightforward consequences of the following
description of the category in the affine case.

